\begin{theorem}[Separating internal and crossing incidence products]%
    \label{thm:peps_graph_open_incidence_product}
    \lean{TNLean.PEPS.prod_region_incident_eq_internal_boundary}
    \leanok
    \uses{def:peps_region_half_edge_label_equiv}
    For a region $R$ in a finite simple graph, a product over its incidences
    equals the product of both endpoint factors on each internal edge,
    multiplied by the product of the single retained endpoint factor on each
    crossing edge. This is the regional incidence decomposition underlying
    \cite[Sections~6.2 and~7]{Schuch2010PEPS}.
\end{theorem}

\begin{proof}\leanok
    Extend incidence factors by $1$ outside $R$, regroup the ambient incidence
    product by edges, and separate internal edges from crossing edges.
\end{proof}

\begin{definition}[Coherent open weighted bond coefficients]%
    \label{def:peps_graph_open_averaging_bond_state}
    \lean{TNLean.PEPS.graphOpenAveragingBondState}
    \leanok
    \uses{def:peps_region_half_edge_label_equiv}
    Let $G$ be finite, $U:G\to\operatorname{Mat}_X(\mathbb C)$ a
    representation, $W$ a virtual weight, and $R$ a region of a finite simple
    graph with ordered endpoints. Write $\kappa=|G|$, and prescribe crossing
    virtual labels $\theta_e$. The physical labels are
    $(\beta_{\partial},\beta_t,\beta_h)$ in crossing, internal-tail, and
    internal-head coordinates. For $q:R\to G$, set
    \begin{align}
        B_e(q) & =
        \begin{cases}
            (WU(q_t^{-1}))_{\theta_e,\beta_{\partial,e}}, & t\in R, \\
            (U(q_h)W)_{\beta_{\partial,e},\theta_e},      & h\in R.
        \end{cases}
    \end{align}
    Define
    \begin{align}
        H_R(\theta,\beta)
        =\kappa^{-|R|}\sum_{q:R\to G}
        \left(\prod_{e\text{ internal}}
        (W^2U(q_hq_t^{-1}))_{\beta_{h,e},\beta_{t,e}}\right)
        \left(\prod_{e\text{ crossing}}B_e(q)\right).
    \end{align}
    These are the regional endpoint factors from the contraction calculation
    in \cite[Sections~6.2 and~7]{Schuch2010PEPS}; the site average is normalized
    as described in \cite{gap:rmp_peps_quantum_double_g_isometry}.
\end{definition}

\begin{theorem}[Actual open weighted averaging contraction]%
    \label{thm:peps_graph_open_averaging_bond_state}
    \lean{TNLean.PEPS.graphOpenRegionNetwork_dressedAveragingSite_coherent,
        TNLean.PEPS.graphOpenBondRegrouping_dressedAveragingSite,
        TNLean.PEPS.graphOpenRegionNetwork_averagingSite_coherent,
        TNLean.PEPS.openRegionWeight_groupBondTensor_dressedAveragingSite}
    \leanok
    \uses{def:peps_graph_open_region_contraction,
        thm:peps_graph_open_region_contraction,
        thm:peps_graph_open_incidence_product,
        def:peps_graph_open_averaging_bond_state,
        lem:peps_region_half_edge_label_endpoints,
        thm:peps_monoid_hom_commuting_weight}
    Suppose $W$ commutes with every $U(g)$. The actual open contraction of
    the averaging sites dressed by $W$, evaluated at a physical incidence
    configuration $\alpha$, equals $H_R(\theta,\beta)$, where $\beta$ is
    the incidence regrouping of $\alpha$. Thus regrouping by the coordinate
    bijection gives exactly the displayed bond formula. The case $W=I$
    gives the ordinary averaging-site coefficients. If the local physical
    incidence configurations admit uniform finite enumerations, the same
    formula is exactly the numbered open-region coefficient of those
    enumerated site tensors. The crossing labels are arbitrary throughout.
    These finite-simple-graph coefficient identities do not assert an
    equivalence of all open-boundary state spaces.
\end{theorem}

\begin{proof}\leanok
    Fix the crossing labels and sum only over internal virtual labels.
    Expand the independent site averages and regroup their incidence products.
    Each internal label is summed once, giving
    $(U(q_h)W)(WU(q_t^{-1}))=W^2U(q_hq_t^{-1})$.
    Crossing labels are fixed, so their single endpoint factors remain.
    The incidence bijection reads the physical endpoint labels literally.
    Finally, finite enumerations reindex the same coefficient sum.
\end{proof}
